\section{Testovanie odporúčacieho systému}

\subsection{Cieľ testovania}

Cieľom testovania bolo overiť správnosť implementácie odporúčacieho algoritmu
a jeho citlivosť na zmenu váhového vektora. Keďže systém pracuje so simulovanými
dátami, nebolo možné merať presnosť odporúčania konkrétnych udalostí — namiesto
toho sme zvolili metriku Precision@5 na úrovni hlavných kategórií.

Po zvážení výsledkov úvodných experimentov (Príloha~\ref{appendix:experiments})
sme dospeli k záveru, že na kvalitu personalizácie majú reálny vplyv iba
\textbf{hlavná kategória a podkategória} — preto sme ďalšie testovanie zamerali
výlučne na tieto dva faktory.

\subsection{Experiment 1 — Overenie správnosti pre realistické profily}

Prvý experiment overoval, či systém správne odporúča udalosti zodpovedajúce
dominantnej kategórii používateľa. Vytvorili sme 5~používateľských profilov
s realistickou históriou:

\begin{itemize}
    \item \textbf{Sporťák} — 6$\times$ Sports, 1$\times$ Art \& Culture, 1$\times$ Food \& Drinks
    \item \textbf{Hudobník} — 5$\times$ Music \& Entertainment, 1$\times$ Sports
    \item \textbf{Umelec} — 4$\times$ Art \& Culture, 1$\times$ Music \& Entertainment
    \item \textbf{Technik} — 3$\times$ Business \& Tech, 1$\times$ Education
    \item \textbf{Foodie} — 4$\times$ Food \& Drinks, 1$\times$ Community \& Social
\end{itemize}

Systém vyhodnotí odporúčanie ako úspešné (\textit{PASS}), ak aspoň 3 z TOP~5
odporúčaní patria do dominantnej kategórie používateľa. Pre každý profil sme
spustili algoritmus s 10~konfiguráciami váh, kde sme postupne menili pomer medzi
hlavnou kategóriou a podkategóriou od 100\,\%/0\,\% až po 0\,\%/100\,\%:

\begin{table}[h]
\centering
\begin{tabular}{|l|c|c|}
\hline
\textbf{Konfigurácia} & \textbf{Hlavná kategória} & \textbf{Podkategória} \\
\hline
W01 & 100\,\% & 0\,\%   \\
W02 & 90\,\%  & 10\,\%  \\
W03 & 80\,\%  & 20\,\%  \\
W04 & 70\,\%  & 30\,\%  \\
W05 & 60\,\%  & 40\,\%  \\
W06 & 50\,\%  & 50\,\%  \\
W07 & 40\,\%  & 60\,\%  \\
W08 & 30\,\%  & 70\,\%  \\
W09 & 20\,\%  & 80\,\%  \\
W10 & 0\,\%   & 100\,\% \\
\hline
\end{tabular}
\caption{Konfigurácie váh pre Experiment 1}
\label{tab:exp1configs}
\end{table}

Systém úspešne prešiel všetkými 10~konfiguráciami pre všetkých 5~používateľov.
Jediný rozdiel sa prejavil pri W10, kde sa na 5.~pozícii začali objavovať udalosti
z iných kategórií — čo potvrdzuje, že hlavná kategória stabilizuje odporúčania a
zabraňuje úniku do nesúvisiacich kategórií.

\subsection{Experiment 2 — Testovanie citlivosti}

V rámci testovania citlivosti sme skúmali, ako zmena histórie používateľa ovplyvňuje
výsledky odporúčacieho systému. Pre každého používateľa sme najprv získali referenčný
zoznam TOP~5 odporúčaní na základe plnej histórie. Následne sme postupne skrývali
každý navštívený event a sledovali, či sa TOP~5 zmenil oproti referenčnému zoznamu.
Zmenu TOP~5 po skrytí eventu považujeme za pozitívny signál — systém reaguje na
aktuálnu históriu a prispôsobuje odporúčania jej zmenám. Testovanie prebehlo na
50~používateľoch naprieč rovnakými 10~konfiguráciami váh.

\begin{table}[h]
\centering
\begin{tabular}{|l|c|c|c|}
\hline
\textbf{Konfigurácia} & \textbf{Hlavná kat.} & \textbf{Podkategória} & \textbf{AVG citlivosť} \\
\hline
T01 & 100\,\% & 0\,\%   & 14\,\% \\
T02 & 90\,\%  & 10\,\%  & 35\,\% \\
T03 & 80\,\%  & 20\,\%  & 35\,\% \\
T04 & 70\,\%  & 30\,\%  & 36\,\% \\
T05 & 60\,\%  & 40\,\%  & 36\,\% \\
T06 & 50\,\%  & 50\,\%  & 40\,\% \\
T07 & 40\,\%  & 60\,\%  & 35\,\% \\
T08 & 30\,\%  & 70\,\%  & 35\,\% \\
T09 & 20\,\%  & 80\,\%  & 34\,\% \\
T10 & 0\,\%   & 100\,\% & 34\,\% \\
\hline
\end{tabular}
\caption{Výsledky testovania citlivosti — podiel prípadov, kde sa TOP~5 zmenil po skrytí eventu}
\label{tab:sensitivity}
\end{table}

Konfigurácia T01 s váhou 100\,\% na hlavnú kategóriu dosiahla najnižšiu citlivosť
14\,\%, čo znamená, že systém reagoval na zmenu histórie len v 14\,\% prípadov.
Naopak, konfigurácia T06 s rovnomerným rozdelením 50\,\%/50\,\% dosiahla najvyššiu
citlivosť 40\,\%. Konfigurácia T10 so 100\,\% váhou na podkategóriu dosiahla len
34\,\%, teda nižšiu citlivosť ako T06. Tento výsledok naznačuje, že kombinácia
oboch faktorov spôsobuje väčšiu reakciu systému na zmeny histórie ako ktorýkoľvek
faktor samostatne.

\subsection{Obmedzenia metodiky}

Leave-one-out metodika má na syntetických dátach niekoľko prirodzených obmedzení.
Používatelia s menej ako 5~navštívenými udalosťami predstavujú cold start scenár,
kde po skrytí jednej udalosti zostáva história príliš malá na spoľahlivú
personalizáciu — ich výsledky by mali byť interpretované samostatne. Pri
diverzifikovanej histórii môže odstránenie jednej udalosti zmeniť dominantnú
kategóriu profilu, pričom systém správne odporučí eventy zodpovedajúce aktuálnej
histórii, ale metrika to hodnotí ako miss — nejde o chybu algoritmu, ale o
obmedzenie metodiky. V neposlednom rade, ak používateľ navštívil určitú kategóriu
iba raz, po jej skrytí klesne hodnota $C_{main}$ na nulu a algoritmus nemá
matematický základ pre odporučenie podobného obsahu. Ani tento jav nepovažujeme
za chybu, ale za prirodzenú vlastnosť Content-based filtra vyžadujúceho aspoň
minimálnu stopu v histórii.

\subsection{Záver testovania}

Testovanie potvrdilo, že algoritmus je deterministický a citlivý na zmenu váhového
vektora. Precision systému je priamo úmerná koncentrácii histórie používateľa —
Content-based filtrovanie funguje najlepšie pre používateľov s jasne definovanými
záujmami. Na základe výsledkov Experimentu~2 sme zvolili konfiguráciu T06
(50\,\%/50\,\%) ako optimálnu pre produkčné nasadenie, keďže kombinácia hlavnej
kategórie a podkategórie v rovnakom pomere zabezpečuje najlepšiu reakciu systému
na zmeny histórie používateľa. Pre komplexnú evalváciu kvality personalizácie by
bolo potrebné testovanie na reálnych používateľských dátach po nasadení aplikácie.
